\documentclass[a4paper,twoside]{article}

\usepackage{morewrites}

\usepackage[svgnames,dvipsnames,x11names]{xcolor}
\usepackage{graphicx}
\usepackage{texsmith-lists}
\usepackage[english]{babel}
\usepackage[colorlinks=true, linkcolor=NavyBlue, urlcolor=NavyBlue, citecolor=NavyBlue]{hyperref}
\usepackage[a4paper]{geometry}
\usepackage{ts-fonts}
\usepackage{float}
\usepackage{seqsplit}
\usepackage{ts-glossary}

\usepackage{lastpage}
\usepackage{refcount}
\makeatletter
\def\ps@plain{
\let\@oddhead\@empty
\let\@evenhead\@empty
\def\@oddfoot{
\hfil
\ifnum\getpagerefnumber{LastPage}>1\relax
\thepage
\fi
\hfil}
\let\@evenfoot\@oddfoot
}
\makeatother

\title{Abbreviations}
\author{}\date{}
\begin{document}
\maketitle
In modern chemical analysis, techniques like \tsacr{NMR}, \tsacr{FTIR}, and \tsacr{GCMS} are essential for identifying compounds, while methods such as \tsacr{HPLC} and \tsacr{ICPOES} provide high-precision quantification of analytes across a wide range of matrices. Researchers often rely on computational tools like \tsacr{DFT} and Molecular Dynamics to model reaction pathways and molecular behavior, supported by data from X-ray Diffraction for crystalline structure determination.

\printglossary[type=\acronymtype]
\end{document}
